\documentclass[a4paper,12pt]{article}

\usepackage[french]{babel}
\usepackage[T1]{fontenc} 
\usepackage[latin1]{inputenc}
%\usepackage{amsmath}

\author{R�mi Synave, Stefka Gueorguieva, Pascal Desbarats}
\title{Manuel de la biblioth�que OVERLAP}
\date{}

\begin{document}
\maketitle

\section{Utilisation}
Biblioth�que o� sont impl�ment�es les fonctions permettant de v�rifier l'intersection de Bounding Box ou de calculer le taux de recouvrement de deux maillages.


\section{Fonctions}


\textbullet \texttt{int intersection\_axis\_aligned\_bounding\_box\_with\_axis\_aligned\_bounding\_box\_list(point3d min, point3d max, point3d *listmin, point3d *listmax, int size)}\\
Test d'intersection entre une Axis Aligned Bounding Box (AABB), \texttt{min} et \texttt{max}, et un ensemble d'AABB, \texttt{listmin} et \texttt{listmax}.\\
~\\
\underline{Param�tres et type de retour :}\\
\texttt{min} : point minimal de la boundingbox � tester (xmin,ymin,zmin).\\
\texttt{max} : point maximal de la boundingbox � tester (xmax,ymax,zmax).\\
\texttt{listmin} : liste de points minimaux des boundingbox (xmin,ymin,zmin).\\
\texttt{listmax} : liste de points maximaux des boundingbox (xmax,ymax,zmax).\\
\texttt{size} : nombre de boundingbox dans la liste.\\
\texttt{retour} : 1 s'il y a instersection, 0 sinon.\\


\textbullet \texttt{int intersection\_bounding\_ball\_with\_bounding\_ball\_list(point3d c, double rayon, point3d *listcentre, double *listrayon, int nbelt, double alpha)}\\
Test d'intersection d'une bounding ball d�finie par son centre \texttt{c} et son rayon \texttt{rayon} avec une liste de bounding ball d�finie par leur centre \texttt{listcentre} et leur rayon \texttt{listrayon}.\\
\textbf{Formule :}\\
Il y a intersection entre la bounding ball de centre $c1$ et de rayon $r1$ et une bounding ball de centre $c2$ et de rayon $r2$ si : $||c2-c1||<||r2-r1|| \times alpha$.\\
~\\
\underline{Param�tres et type de retour :}\\
\texttt{c} : centre de la bounding sphere � tester.\\
\texttt{rayon} : rayon de la  bounding sphere � tester.\\
\texttt{listcentre} : liste de centre des bounding sphere.\\
\texttt{listrayon} : liste des rayons des bounding sphere.\\
\texttt{nbelt} : taille de la liste de bounding sphere.\\
\texttt{alpha} : param�tre de "rapprochement" (voir formule).\\
\texttt{retour} : 1 s'il y a instersection, 0 sinon.\\



\textbullet \texttt{double vf\_model\_bounding\_ball\_ritter\_compute\_overlap(vf\_model *base, vf\_model *m)}\\
Calcul du taux de recouvrement de deux maillages en utilisant les bounding ball de Ritter.\\
\textbf{Algorithme :}\\
1) Calculer les bounding ball de Ritter de toutes les faces des mod�les \texttt{base} et \texttt{m}.\\
2) Pour toutes les faces de \texttt{m}, si la bounding ball de cette face s'intersecte avec l'une des bounding ball des faces de \texttt{base} alors marquer la face de \texttt{m}.\\
3) Pour toutes les faces marqu�es de \texttt{m}, ajouter les sommets sans redondance � une liste \texttt{L}.\\
4) Le taux de recouvrement = $\frac {||L||} {||m||} $.\\
o� ||\texttt{L}|| est le nombre d'�l�ments de la liste \texttt{L} et ||\texttt{m}| est le nombre de sommets du maillage \texttt{m}.\\
~\\
\underline{Param�tres et type de retour :}\\
\texttt{base} : premier modele.\\
\texttt{m} : second modele.\\
\texttt{retour} : taux de recouvrement de \texttt{m} sur \texttt{base}.\\



\textbullet \texttt{double vf\_model\_axis\_aligned\_bounding\_box\_compute\_overlap(vf\_model *base, vf\_model *m)}\\
Calcul du taux de recouvrement de deux maillages en utilisant les AABB.\\
\textbf{Algorithme :}\\
1) Calculer les AABB de toutes les faces des mod�les \texttt{base} et \texttt{m}.\\
2) Pour toutes les faces de \texttt{m}, si l'AABB de cette face s'intersecte avec l'une des AABB des faces de \texttt{base} alors marquer la face de \texttt{m}.\\
3) Pour toutes les faces marqu�es de \texttt{m}, ajouter les sommets sans redondance � une liste \texttt{L}.\\
4) Le taux de recouvrement = $\frac {||L||} {||m||} $.\\
o� ||\texttt{L}|| est le nombre d'�l�ments de la liste \texttt{L} et ||\texttt{m}| est le nombre de sommets du maillage \texttt{m}.\\
~\\
\underline{Param�tres et type de retour :}\\
\texttt{base} : premier modele.\\
\texttt{m} : second modele.\\
\texttt{retour} : taux de recouvrement de \texttt{m} sur \texttt{base}.\\


\textbullet \texttt{void vf\_model\_bounding\_ball\_ritter\_find\_face\_overlapping(vf\_model *base, vf\_model *m, int **list, int *size, double alpha)}\\
Liste les faces de \texttt{m} s'intersectant avec les face de \texttt{base} en utilisant les bounding ball de Ritter. L'indes des faces est ajout�es � la liste \texttt{list}.\\
\textbf{Algorithme :}\\
1) Calculer les bounding ball de Ritter de toutes les faces des mod�les \texttt{base} et \texttt{m}.\\
2) Pour toutes les faces de \texttt{m}, si la bounding ball de cette face s'intersecte avec l'une des bounding ball des faces de \texttt{base} alors ajouter la face � la liste \texttt{list}.\\
~\\
\underline{Param�tres et type de retour :}\\
\texttt{base} : premier modele.\\
\texttt{m} : second modele.\\
\texttt{list} : liste des num�ros des faces de \texttt{m} s'intersectant avec \texttt{base}.\\
\texttt{size} : taille de la liste.\\
\texttt{alpha} : param�tre de rapprochement (voir formule de la fonction intersection\_bounding\_ball\_with\_bounding\_ball\_list).\\
\texttt{retour} : aucun.\\


\textbullet \texttt{void vf\_model\_axis\_aligned\_bounding\_box\_find\_face\_overlapping(vf\_model *base, vf\_model *m, int **listface, int *size)}\\
Liste les faces de \texttt{m} s'intersectant avec les face de \texttt{base} en utilisant les AABB. L'indes des faces est ajout�es � la liste \texttt{list}.\\
\textbf{Algorithme :}\\
1) Calculer les AABB de toutes les faces des mod�les \texttt{base} et \texttt{m}.\\
2) Pour toutes les faces de \texttt{m}, si l'AABB de cette face s'intersecte avec l'une des AABB des faces de \texttt{base} alors ajouter la face � la liste \texttt{list}.\\
~\\
\underline{Param�tres et type de retour :}\\
\texttt{base} : premier modele.\\
\texttt{m} : second modele.\\
\texttt{list} : liste des num�ros des faces de \texttt{m} s'intersectant avec \texttt{base}.\\
\texttt{size} : taille de la liste.\\
\texttt{retour} : aucun.\\


\end{document}
